\documentclass[letterpaper,11pt]{article}

\usepackage[empty]{fullpage}
\usepackage[hidelinks]{hyperref}
\usepackage{fontawesome5}
\usepackage{enumitem}
\usepackage{parskip}

% Adjust margins
\addtolength{\oddsidemargin}{-0.4in}
\addtolength{\evensidemargin}{-0.4in}
\addtolength{\textwidth}{0.8in}
\addtolength{\topmargin}{-0.5in}
\addtolength{\textheight}{1.0in}

\ifdefined\pdfgentounicode
  \input{glyphtounicode}
  \pdfgentounicode=1
\fi

\begin{document}
\begin{center}
    {\Huge \scshape Juan Becerra} \\ \vspace{4pt}
    \textbf{\large Software Engineer -- Engine, Gameplay \& Systems} \\ \vspace{4pt}
    \small Seattle, WA ~ $|$ ~
    \raisebox{-0.1\height}\faPhone\ (208) 550-9191 ~ $|$ ~
    \href{mailto:becerrajuan007@gmail.com}{\raisebox{-0.2\height}\faEnvelope\ \underline{becerrajuan007@gmail.com}} ~ $|$ ~
    \href{https://hoowan.dev}{\raisebox{-0.2\height}\faGlobe\ \underline{hoowan.dev}} \\ \vspace{2pt}
    \href{https://linkedin.com/in/hoowan/}{\raisebox{-0.2\height}\faLinkedin\ \underline{linkedin.com/in/hoowan}} ~ $|$ ~
    \href{https://github.com/hoowan-dev}{\raisebox{-0.2\height}\faGithub\ \underline{github.com/hoowan-dev}}
\end{center}

\vspace{4pt}
\hrule
\vspace{14pt}

\noindent November 3, 2026

\vspace{8pt}
\noindent \textbf{Worldscape Hiring Team} \\
Worldscape \\
Redmond, WA

\vspace{8pt}
\noindent \textbf{Re: Game Engine Engineer, Platform}

\vspace{8pt}
\noindent Dear Worldscape Hiring Team,

I am writing to express my strong interest in the \textbf{Game Engine Engineer, Platform} position at
Worldscape. Over the past five years at Electronic Arts, as a Gameplay \& Animation Software Engineer at
\textbf{Maxis} on \emph{Project Rene} (\emph{The Sims}) and now as a Tools \& Tech-Art Software Engineer at
\textbf{EA Lifestyle}, I have built real-time simulation systems in modern C++ under strict frame budgets,
along with the tools that help designers and QA iterate on them. Worldscape's focus on physics-based 3D
simulation for digital twins is an exciting place to apply that experience.

At \textbf{EA Maxis Studios}, I architected a multi-threaded, asynchronous path-planning and navigation API in
Unreal Engine 5 and C++ that delivered deterministic crowd routing and real-time obstacle avoidance within a
frame budget of under 1.5ms. I profiled, benchmarked, and optimized compute, memory, and frame budgets across
character evaluation graphs and spatial query sweeps to eliminate hitches in dense crowd simulations. I also
implemented core locomotion, character physics, state machines, event systems, and runtime animation,
including motion matching, distance matching, and custom IK solvers, which is exactly the kind of
physics-to-animation integration your team describes. To keep iteration fast and stable, I partnered with
designers, technical animators, and QA to build in-editor tuning tools and visual debuggers.

At \textbf{EA Lifestyle}, I engineered a high-throughput Unreal Engine C++ plugin that captures runtime character
trajectory, pose, and velocity state for motion-matching databases, and built a Maya/Python authoring tool
adopted across multiple studios that cut redundant keyframing work by over 70\%. Earlier, as a Senior Games
Test Engineer at \textbf{Lionbridge}, I led regression testing on live-service AAA titles, which taught me how
much stability and clear repro steps matter once content ships.

Outside of studio work, I built a custom C++ engine (\emph{Hoowan Game Engine}) with collision detection and
correction, an EnTT-based entity-component system, an OpenGL renderer, and a desktop editor. I also bring
hands-on experience with broad-phase and continuous collision detection, integration schemes and constraint
solvers, Unreal Chaos and Unity DOTS, data-oriented design, GPU compute, and building test suites to track
down crashes and regressions.

Based in Seattle, I would be glad to work on-site in Redmond and help Worldscape build fast, robust physics
for real-time simulation. Thank you for your time and consideration. I look forward to speaking with you.

\vspace{10pt}
\noindent Sincerely,

\vspace{14pt}
\noindent \textbf{Juan Becerra}

\end{document}
